% =====================================================================
%  Lecture 26 -- Future 6G Technologies                   (English edition)
% =====================================================================
\documentclass[11pt,a4paper]{article}
\providecommand\ACNroot{../..}
\input{\ACNroot/templates/lecture-en}

\ACNset{number}{26}
\ACNset{title}{Future 6G Technologies}
\ACNset{shorttitle}{6G technologies}
\ACNset{titleother}{فناوری‌های آیندهٔ نسل ششم}
\ACNset{part}{Part V -- Planes and outlook (Lectures 23--26)}
\ACNset{time}{6--8 hours (3 h reading, 3--5 h exercises)}
\ACNset{level}{Advanced}
\ACNset{prereq}{Lectures 14, 23 and 25; Shannon capacity and free-space propagation}
\ACNset{gen}{6}

\begin{document}
\ACNcover

% ---------------------------------------------------------------------
\section{Learning objectives}
After studying this lecture you should be able to:
\begin{objectives}
  \item \textbf{Describe} the IMT-2030 framework: its usage scenarios and overarching aspects.
  \item \textbf{Explain} the main candidate technologies: sub-THz spectrum, cell-free massive MIMO, reconfigurable intelligent surfaces,
        integrated sensing and communication, non-terrestrial networks and AI-native design.
  \item \textbf{Compute} free-space path loss at high frequencies and the capacity gain of wider bandwidth under a power limit.
  \item \textbf{Estimate} RIS gain, sensing resolution and satellite propagation delay.
  \item \textbf{Assess} which claims about 6G are established and which are still research.
\end{objectives}

% ---------------------------------------------------------------------
\section{Prerequisite knowledge}
You need the 5G requirements and NR design (Lecture~14), the user-plane budgets of Lecture~23 and the evolution trends of Lecture~25, as well as
the Shannon capacity formula and the free-space propagation (Friis) equation.

\paragraph{Key terms} The vocabulary of this lecture with its Persian equivalents:
\begin{keyterms}
IMT-2030 & \fa{IMT-2030} & ITU framework for the sixth generation (Recommendation M.2160). \\
Sub-THz & \fa{زیرتراهرتز} & Carrier frequencies of roughly 100--300\,GHz. \\
Cell-free massive MIMO & \fa{MIMO انبوه بدون سلول} & Many distributed access points jointly serving all users. \\
RIS & \fa{سطح هوشمند بازپیکربندی‌پذیر} & Surface of passive elements with controllable phase that shapes reflections. \\
ISAC & \fa{حسگری و ارتباط یکپارچه} & Using the same signals and hardware for radar sensing and communication. \\
NTN & \fa{شبکهٔ غیرزمینی} & Satellites and high-altitude platforms as part of the mobile network. \\
AI-native & \fa{هوش‌مصنوعی‌بنیاد} & Designed with learned components and data pipelines from the start. \\
Range resolution & \fa{تفکیک‌پذیری فاصله} & Smallest distance difference between two distinguishable targets. \\
Power-limited regime & \fa{رژیم محدود به توان} & Where more bandwidth no longer increases capacity much. \\
Energy per bit & \fa{انرژی به ازای بیت} & Consumed energy divided by delivered bits (J/bit). \\
\end{keyterms}

% ---------------------------------------------------------------------
\section{Introduction and motivation}
Each generation has been defined about ten years after the previous one, and 6G follows the same rhythm: the ITU published the IMT-2030
framework in 2023 \cite{itu2023m2160}, and 3GPP began 6G study work in Release~20, with normative specifications expected in later
releases. As Lecture~25 showed, the next steps in capacity need much wider bandwidths, the cost per bit must keep falling, and large parts of
the world are still unconnected. At the same time, the network is expected to do more than move bits: sense its environment, host AI and be
sustainable and resilient \cite{saad2020vision}.

Much of 6G is still research. This lecture therefore separates the agreed framework from candidate technologies, and gives simple physical
models that let you judge claims about them yourself.

% ---------------------------------------------------------------------
\section{Fundamental concepts}

\subsection{The IMT-2030 framework}
Recommendation ITU-R M.2160 \cite{itu2023m2160} defines six usage scenarios (\cref{fig:usage}). Three extend the 5G triangle: \emph{immersive
communication} (from eMBB), \emph{hyper-reliable and low-latency communication} (HRLLC, from URLLC) and \emph{massive communication} (from mMTC).
Three are new: \emph{ubiquitous connectivity}, \emph{AI and communication} and \emph{integrated sensing and communication}. Four overarching
aspects apply to all: sustainability, connecting the unconnected, ubiquitous intelligence, and security and resilience. Quantitative
requirements are defined in a later ITU step, so numbers such as ``1\,Tbit/s'' found in the literature are research targets, not
requirements.

\begin{acnfigure}
  \resizebox{\linewidth}{!}{\acnfig{lecture-26/usage-scenarios}}
  \caption{IMT-2030 usage scenarios and overarching aspects (after ITU-R M.2160).}
  \label{fig:usage}
\end{acnfigure}

\subsection{Candidate technologies}
\begin{itemize}
  \item \textbf{New spectrum}: upper mid-band (around 7--24\,GHz) for wide-area capacity and sub-THz bands for very high rates at short range
        \cite{rappaport2019above100}.
  \item \textbf{Cell-free massive MIMO}: many access points coherently serving users without cell boundaries, giving uniform service
        \cite{ngo2017cellfree}.
  \item \textbf{Reconfigurable intelligent surfaces}: controllable reflectors that create paths around obstacles at low power \cite{wu2019irs}.
  \item \textbf{Integrated sensing and communication}: the radio as a radar for positioning, mapping and gesture or object detection
        \cite{liu2022isac}.
  \item \textbf{Non-terrestrial networks}: LEO satellites and high-altitude platforms integrated with the terrestrial system \cite{tr38811}.
  \item \textbf{AI-native air interface and network}: learned receivers, beam management and automated operation.
\end{itemize}

% ---------------------------------------------------------------------
\section{Detailed technical explanation}

\subsection{Sub-THz communication}
Above 100\,GHz contiguous bandwidths of tens of gigahertz are available, but path loss grows with frequency for antennas of fixed gain, molecular
absorption appears at specific frequencies, and blockage by walls, bodies and even foliage is severe. Very large arrays with narrow beams
recover the loss (\cref{sec:fspl}), which makes beam search and tracking the central problem. Hardware is also hard: power amplifiers are
inefficient, phase noise grows with frequency, and sampling tens of gigahertz needs low-resolution converters.

\subsection{Cell-free networks}
In a cellular network users at cell edges suffer from interference and weak signals. In a cell-free network a central processing unit
coordinates many distributed access points that all serve all users, with local precoding and fronthaul to the central unit. The result is
more uniform rates; the costs are fronthaul and synchronisation -- an extension of the CU/DU split and coordination ideas of Lecture~23.

\subsection{Reconfigurable intelligent surfaces}
An RIS is a panel of many sub-wavelength elements whose reflection phase can be set electronically. Properly configured, it adds a strong
reflected path, for example around a blocking building. Its received power grows with the square of the number of elements
(\cref{sec:ris}), but it suffers the product of two path losses (transmitter to surface, surface to receiver), so it helps most near the
transmitter or the receiver.

\subsection{Integrated sensing and communication}
Wide bandwidths and large arrays are exactly what radars need. ISAC uses communication waveforms or shared hardware to measure range, angle
and velocity of objects. Applications include positioning without dedicated infrastructure, traffic monitoring and environment mapping that in
turn helps beam management. It also raises privacy questions, since the network could sense people who carry no device.

\subsection{Non-terrestrial networks}
Release~17 introduced NR over satellite (NTN) with adaptations for large delays and Doppler shifts. LEO constellations offer global coverage
with tens of milliseconds of round-trip delay (\cref{sec:leo}); integrated with terrestrial networks they address ``connecting the
unconnected'' and resilience after disasters.

\subsection{AI-native networks}
5G already uses data analytics (NWDAF) and AI in RAN optimisation. 6G aims to make learned components part of the design: AI-based channel
state feedback, receivers and beam prediction, and network automation with closed loops. Open questions are standardisation of models,
generalisation to unseen conditions, energy cost of the AI itself and explainability.

% ---------------------------------------------------------------------
\section{Mathematical models}

\subsection{Free-space path loss}\label{sec:fspl}
\paragraph{Model} Between isotropic antennas at distance $d$ and wavelength $\lambda = c/f$:
\begin{equation}\label{eq:fspl}
  L_{\mathrm{FS}} = 20\log_{10}\!\left(\frac{4\pi d f}{c}\right) \ \mathrm{dB}.
\end{equation}
With effective apertures $A_t$, $A_r$ the Friis equation becomes $P_r = P_t A_t A_r/(\lambda^2 d^2)$, which \emph{grows} with $f^2$.
\paragraph{Example} At 100\,m: 83.3\,dB at 3.5\,GHz, 101.4\,dB at 28\,GHz and 115.4\,dB at 140\,GHz (\cref{fig:fspl}) -- 32\,dB more at 140\,GHz than at 3.5\,GHz.
Arrays of the same physical size at both ends gain $20\log_{10}(140/3.5) = 32$\,dB each, so the link budget can actually improve, at the price of
very narrow beams.

\begin{acnfigure}
  \acnfig{lecture-26/fspl}
  \caption{Free-space path loss (\cref{eq:fspl}).}
  \label{fig:fspl}
\end{acnfigure}

\subsection{Bandwidth under a power limit}\label{sec:bw}
\paragraph{Model} With received power $P$ and noise density $N_0$:
\begin{equation}\label{eq:cap}
  C(B) = B \log_2\!\left(1 + \frac{P}{N_0 B}\right) \ \xrightarrow{B\to\infty}\ \frac{P}{N_0 \ln 2} .
\end{equation}
\paragraph{Example} If the SNR is 20\,dB over 1\,GHz: $C = 6.66$\,Gbit/s; with 10\,GHz, 34.6\,Gbit/s; with 100\,GHz, 100\,Gbit/s; the limit is 144\,Gbit/s
(\cref{fig:bw}). Wider bandwidth helps greatly only while the SNR stays high -- hence the need for array gain.

\begin{acnfigure}
  \acnfig{lecture-26/bandwidth-capacity}
  \caption{Capacity versus bandwidth at fixed received power (\cref{eq:cap}).}
  \label{fig:bw}
\end{acnfigure}

\subsection{RIS gain}\label{sec:ris}
\paragraph{Model} With $N$ elements whose phases align the reflected signals at the receiver, the amplitudes add coherently:
\begin{equation}\label{eq:ris}
  P_r \propto N^2 , \qquad G_{\mathrm{RIS}} = 20\log_{10} N \ \mathrm{dB}
\end{equation}
relative to a single element.
\paragraph{Example} $N = 1024$: 60.2\,dB over one element (random phases would give only $10\log_{10}N = 30.1$\,dB). The gain must overcome the product
of the two path losses.

\subsection{Sensing resolution}\label{sec:isac}
\paragraph{Model} A signal of bandwidth $B$ separates two targets in range if
\begin{equation}\label{eq:range}
  \Delta r = \frac{c}{2B} .
\end{equation}
\paragraph{Example} $B = 1$\,GHz: 15\,cm; $B = 10$\,GHz: 1.5\,cm. Sub-THz bandwidths thus allow centimetre-level sensing.

\subsection{Satellite delay}\label{sec:leo}
\paragraph{Model} For a satellite at altitude $h$, Earth radius $R$ and elevation angle $\varepsilon$, the slant range is
\begin{equation}\label{eq:slant}
  d(\varepsilon) = \sqrt{(R + h)^2 - R^2\cos^2\varepsilon} - R\sin\varepsilon ,
\end{equation}
and the one-way delay $d/c$.
\paragraph{Example} $h = 600$\,km: 600\,km and 2.0\,ms at zenith, 1076\,km and 3.6\,ms at $30^\circ$, 1932\,km and 6.4\,ms at $10^\circ$ (\cref{fig:leo}). A
geostationary satellite (35\,786\,km) needs at least 119\,ms.

\begin{acnfigure}
  \acnfig{lecture-26/leo-delay}
  \caption{One-way propagation delay to a LEO satellite (\cref{eq:slant}).}
  \label{fig:leo}
\end{acnfigure}

% ---------------------------------------------------------------------
\section{Network architecture and diagrams}
A plausible 6G system keeps the service-based core of 5G (Lecture~18) and extends it with sensing and AI functions, integrates satellite and
terrestrial access under one core, and deploys a RAN that mixes macro cells in the upper mid-band, dense sub-THz spots, cell-free clusters
and RIS-assisted coverage. \Cref{fig:usage} gives the service view; \cref{fig:fspl,fig:bw,fig:leo} quantify the physical constraints.

% ---------------------------------------------------------------------
\section{Step-by-step procedures}

\begin{procedurebox}{Checking a sub-THz link budget}\label{proc:budget}
\begin{steps}
  \item Compute the free-space loss at the carrier and distance (\cref{eq:fspl}); add absorption and blockage margins.
  \item Add transmit and receive array gains from the array sizes.
  \item Compute the noise power $N_0B$ including the receiver noise figure, and the SNR.
  \item Compute the capacity (\cref{eq:cap}) and compare with the target rate.
\end{steps}
\end{procedurebox}

\begin{procedurebox}{Judging a 6G claim}\label{proc:claim}
\begin{steps}
  \item Is it a requirement in an ITU or 3GPP document or a research target?
  \item Which physical quantity limits it (bandwidth, SNR, delay, energy)? Check it with a model of this lecture.
  \item What does it cost (spectrum, sites, energy, complexity)?
  \item Which usage scenario of \cref{fig:usage} does it serve?
\end{steps}
\end{procedurebox}

% ---------------------------------------------------------------------
\section{Worked examples}

\begin{examplebox}{Short-range sub-THz link}\label{ex:thz}
\textbf{Problem.} Compute the free-space loss at 300\,GHz over 10\,m.

\textbf{Solution.} $20\log_{10}(4\pi\times10\times3\times10^{11}/3\times10^{8}) = 102.0$\,dB -- about the same as 28\,GHz over 100\,m, which is why sub-THz suits short
ranges.
\end{examplebox}

\begin{examplebox}{More bandwidth}\label{ex:bw}
\textbf{Problem.} With the power of \cref{sec:bw}, what capacity does 50\,GHz give?

\textbf{Solution.} SNR $= 100/50 = 2$; $C = 50\log_2 3 = 79.2$\,Gbit/s: 50 times the bandwidth of the 1\,GHz case gives only 12 times the capacity.
\end{examplebox}

\begin{examplebox}{Satellite round trip}\label{ex:leo}
\textbf{Problem.} A UE and a gateway both see a 600\,km LEO satellite at $30^\circ$. What is the round-trip propagation delay through the satellite (bent pipe)?

\textbf{Solution.} One way UE--satellite--gateway: $2\times3.58 = 7.2$\,ms; round trip 14.3\,ms -- compatible with interactive services, unlike GEO (over 477\,ms).
\end{examplebox}

% ---------------------------------------------------------------------
\section{Technology comparison}
\begin{table}[htbp]
\centering\small
\caption{5G and candidate 6G features.}\label{tab:cmp}
\rowcolors{2}{acnmist}{white}
\begin{tabularx}{\linewidth}{@{}l Y Y@{}}
\toprule
\rowcolor{white} & 5G (IMT-2020) & 6G (IMT-2030, candidates) \\
\midrule
Usage scenarios & eMBB, URLLC, mMTC & Immersive, HRLLC, massive, ubiquitous, AI, sensing \\
Spectrum & Below 7\,GHz and 24--71\,GHz & + upper mid-band and sub-THz \\
Radio architecture & Cells with massive MIMO & Cells + cell-free, RIS \\
Sensing & Positioning & Integrated sensing and communication \\
Non-terrestrial & NTN from Release~17 & Integrated by design \\
AI & Analytics (NWDAF), RAN optimisation & AI-native air interface and operation \\
Status & Deployed & Research and early standardisation \\
\bottomrule
\end{tabularx}
\end{table}

\paragraph{Evolution, not revolution} As every generation before it, 6G will reuse the persistent functions of Lecture~25 and add new
capabilities where physics and economics allow.

% ---------------------------------------------------------------------
\section{Key takeaways}
\begin{takeawaybox}
\begin{itemize}
  \item IMT-2030 defines six usage scenarios and four overarching aspects; numerical requirements come later.
  \item Free-space loss grows as $20\log_{10} f$, but fixed-aperture arrays recover it: beams, not power, are the problem.
  \item Bandwidth helps only while SNR is high: $C \to P/(N_0 \ln 2)$.
  \item RIS gain scales as $N^2$; sensing resolution as $c/(2B)$; LEO delay is a few milliseconds.
  \item Separate requirements from research targets when reading about 6G.
\end{itemize}
\end{takeawaybox}

% ---------------------------------------------------------------------
\section{Exercises}

\begin{exercisebox}{1}{Scenarios}\label{exr:scen}
Name the three IMT-2030 usage scenarios that extend 5G and the three new ones.
\end{exercisebox}

\begin{exercisebox}{1}{Blockage}\label{exr:block}
Give two reasons why sub-THz links are more sensitive to blockage than links at 3.5\,GHz, and one mitigation.
\end{exercisebox}

\begin{exercisebox}{2}{Radar resolution}\label{exr:range}
Compute the range resolution of a 400\,MHz NR carrier used for sensing.
\end{exercisebox}

\begin{exercisebox}{2}{RIS size}\label{exr:ris}
A 16-element RIS gives too weak a reflection. How many elements are needed for 20\,dB more received power?
\end{exercisebox}

\begin{exercisebox}{3}{GEO versus LEO}\label{exr:geo}
Compute the minimum round-trip delay UE--satellite--gateway for a GEO satellite at 35\,786\,km and compare with \cref{ex:leo}. Which services are feasible?
\end{exercisebox}

\begin{exercisebox}{3}{Energy}\label{exr:energy}
Traffic grows 30\,\% per year while the network energy budget stays constant. By how much must energy per bit improve per year and over ten years?
\end{exercisebox}

\subsection*{Solutions}
\addcontentsline{toc}{subsection}{Solutions}

\begin{solutionbox}{\ref{exr:scen}}
Extended: immersive communication, HRLLC, massive communication. New: ubiquitous connectivity, AI and communication, integrated sensing and communication.
\end{solutionbox}

\begin{solutionbox}{\ref{exr:block}}
Shorter wavelength means less diffraction around obstacles and higher penetration loss through materials; narrow beams are easily blocked. Mitigations:
multi-connectivity to several access points, RIS-assisted paths, fast beam switching.
\end{solutionbox}

\begin{solutionbox}{\ref{exr:range}}
$\Delta r = 3\times10^{8}/(2\times4\times10^{8}) = 0.375$\,m.
\end{solutionbox}

\begin{solutionbox}{\ref{exr:ris}}
$(N/16)^2 = 100 \Rightarrow N = 160$ elements.
\end{solutionbox}

\begin{solutionbox}{\ref{exr:geo}}
Each hop at least 35\,786\,km (119.3\,ms): one way 238.6\,ms, round trip 477\,ms, about 33 times the LEO value. Voice is degraded and interactive or
real-time control services are not feasible; broadcasting and bulk data work well.
\end{solutionbox}

\begin{solutionbox}{\ref{exr:energy}}
Energy per bit must fall by a factor 1.3 each year (23\,\%), i.e.\ $1.3^{10} = 13.8$ times in ten years -- hence the emphasis on sleep modes, efficient
hardware and AI-driven energy saving.
\end{solutionbox}

% ---------------------------------------------------------------------
\section{Further reading}
The IMT-2030 framework is Recommendation ITU-R M.2160 \cite{itu2023m2160}. Saad et al.\ \cite{saad2020vision} present a 6G vision; Rappaport et al.\
\cite{rappaport2019above100} discuss frequencies above 100\,GHz; Ngo et al.\ \cite{ngo2017cellfree} introduce cell-free massive MIMO; Wu and Zhang \cite{wu2019irs}
analyse intelligent reflecting surfaces; Liu et al.\ \cite{liu2022isac} survey ISAC; TR~38.811 \cite{tr38811} studies NR for non-terrestrial networks.

\printbibliography[heading=subbibliography,title={\ACNstrRefs}]
\end{document}
